% ===== PRÉAMBULE CANONIQUE — à coller VERBATIM en tête de CHAQUE .tex =====
\documentclass[11pt,a4paper]{article}
\usepackage[utf8]{inputenc}\usepackage[T1]{fontenc}\usepackage{lmodern}
\usepackage[french]{babel}
\usepackage{amsmath,amssymb,mathtools}\usepackage{icomma}
\usepackage[version=4]{mhchem}          % \ce{...} : équations & formules
\usepackage{siunitx}\sisetup{locale=FR} % \qty{}{}, \si{}, \ang{}
\usepackage{chemfig}                    % \chemfig{...} : structures développées
\usepackage{xcolor}\usepackage{colortbl}
\usepackage{tikz}\usepackage{pgfplots}\pgfplotsset{compat=1.18}
\usepackage[most]{tcolorbox}\usepackage{enumitem}\usepackage{array,booktabs}
\usepackage[a4paper,margin=2.2cm]{geometry}
\usetikzlibrary{arrows.meta,calc,angles,quotes,positioning,patterns,backgrounds,shapes.geometric}
\definecolor{primary}{RGB}{222,49,99}\definecolor{primarydark}{RGB}{150,22,55}
\definecolor{defcol}{RGB}{0,114,178}\definecolor{methcol}{RGB}{52,92,148}
\definecolor{excol}{RGB}{0,135,90}\definecolor{attncol}{RGB}{200,40,55}
\tcbset{boxbase/.style={enhanced,breakable,boxrule=0.9pt,arc=2.5pt,
  left=3mm,right=3mm,top=2.3mm,bottom=2.3mm,fonttitle=\bfseries,coltitle=white}}
% --- boîtes TUTORIEL (noms EXACTS, ne pas renommer) ---
\newtcolorbox{cle}[1][]{boxbase,colback=defcol!6,colframe=defcol,
  title={\textsc{À retenir}\ifx\relax#1\relax\relax\else\textmd{ : \itshape #1}\fi}}
\newtcolorbox{methode}[1][]{boxbase,colback=methcol!7,colframe=methcol,sharp corners,
  title={\textsc{Méthode}\ifx\relax#1\relax\relax\else\textmd{ --- \itshape #1}\fi}}
\newtcolorbox{exemplebox}[1][]{boxbase,colback=excol!5,colframe=excol,
  title={\textsc{Exemple}\ifx\relax#1\relax\relax\else\textmd{ : \itshape #1}\fi}}
\newtcolorbox{piege}[1][]{boxbase,colback=attncol!6,colframe=attncol,
  title={\textsc{Piège}\ifx\relax#1\relax\relax\else\textmd{ : \itshape #1}\fi}}
\newtcolorbox{remarque}[1][]{boxbase,colback=black!4,colframe=black!45,
  title={\textsc{Remarque}\ifx\relax#1\relax\relax\else\textmd{ : \itshape #1}\fi}}
% --- boîtes EXERCICE / CORRIGÉ (noms EXACTS, auto-comptées) ---
\newtcolorbox[auto counter]{exo}[1][]{boxbase,colback=primary!4,colframe=primary,
  title={\textsc{Exercice}~\thetcbcounter\ifx\relax#1\relax\relax\else\textmd{ --- \itshape #1}\fi}}
\newtcolorbox[auto counter]{corrige}[1][]{boxbase,colback=excol!4,colframe=excol,
  title={\textsc{Corrigé}~\thetcbcounter\ifx\relax#1\relax\relax\else\textmd{ --- \itshape #1}\fi}}
\newcommand{\R}{\mathbb{R}}\newcommand{\N}{\mathbb{N}}\newcommand{\Z}{\mathbb{Z}}\newcommand{\ds}{\displaystyle}
% (un \usepackage{fancyhdr}+en-tête est possible ; il est ignoré côté web)
% =========================================================================
\begin{document}
\newcommand{\CF}{\ensuremath{\mathrm{C.F.}}}
\begin{corrige}[toutes les charges de \ce{CO3^2-}]
$\CF = N_v - 2\times(\text{doublets libres}) - (\text{liaisons})$.
\begin{enumerate}[label=\alph*)]
  \item \ce{O} doublement lié : $6 - 2\times 2 - 2 = \mathbf{0}$.
  \item \ce{O} simplement lié : $6 - 2\times 3 - 1 = \mathbf{-1}$.
  \item \ce{C} : $4 - 0 - 4 = \mathbf{0}$.
  \item Somme : $0 + 0 + (-1) + (-1) = \mathbf{-2}$ : c'est bien la charge de \ce{CO3^2-}. \checkmark
\end{enumerate}
\end{corrige}

\begin{corrige}[toutes les charges de \ce{NO3-}]
\begin{enumerate}[label=\alph*)]
  \item \ce{N} : $5 - 0 - 4 = \mathbf{+1}$.
  \item \ce{O} doublement lié : $6 - 4 - 2 = \mathbf{0}$.
  \item \ce{O} simplement lié : $6 - 6 - 1 = \mathbf{-1}$.
  \item Somme : $(+1) + 0 + (-1) + (-1) = \mathbf{-1}$ : c'est la charge de \ce{NO3-}. \checkmark
\end{enumerate}
\end{corrige}

\begin{corrige}[une neutre à charges cachées : \ce{CO}]
Triple liaison $\Rightarrow$ $3$ liaisons ; chaque atome porte $1$ doublet libre.
\begin{enumerate}[label=\alph*)]
  \item \ce{C} : $4 - 2\times 1 - 3 = \mathbf{-1}$.
  \item \ce{O} : $6 - 2\times 1 - 3 = \mathbf{+1}$.
  \item Somme : $(-1) + (+1) = \mathbf{0}$. C'est cohérent : \ce{CO} est neutre, mais porte deux
        charges formelles opposées qui se compensent. La charge formelle n'est donc pas la charge
        « réelle » de l'atome.
\end{enumerate}
\end{corrige}

\begin{corrige}[l'ozone \ce{O3}]
\begin{enumerate}[label=\alph*)]
  \item \ce{O} central ($1$ doublet, $3$ liaisons : une double $+$ une simple) :
        $6 - 2 - 3 = \mathbf{+1}$.
  \item \ce{O} doublement lié : $6 - 4 - 2 = \mathbf{0}$.
  \item \ce{O} simplement lié : $6 - 6 - 1 = \mathbf{-1}$.
  \item Somme : $(+1) + 0 + (-1) = \mathbf{0}$ : l'ozone est neutre, avec une charge $+$ sur
        l'atome central et une charge $-$ sur un atome terminal.
\end{enumerate}
\end{corrige}

\begin{corrige}[choisir la meilleure structure : \ce{CO2}]
\begin{enumerate}[label=\alph*)]
  \item Structure (i) \ce{O=C=O} : chaque \ce{O} $= 6-4-2 = 0$ ; \ce{C} $= 4-0-4 = 0$.
        \textbf{Toutes les \CF{} sont nulles.}
  \item Structure (ii) \ce{O#C-O} : \ce{O} triplement lié $= 6-2-3 = +1$ ; \ce{C} $= 4-0-4 = 0$ ;
        \ce{O} simplement lié $= 6-6-1 = -1$. \textbf{Charges $+1$, $0$, $-1$.}
  \item La structure \textbf{(i)} est la meilleure : ses charges formelles sont \emph{minimales}
        (toutes nulles), alors que (ii) fait apparaître des charges $\pm 1$ inutiles.
\end{enumerate}
\end{corrige}

\end{document}
